\documentclass[autodetect-engine,11pt,a4paper]{bxjsarticle}
\usepackage{amsmath,amssymb}
\usepackage{geometry}
\usepackage{bm}
\geometry{margin=10mm}

\title{\Large\textbf{宇宙の流体力学 レポート課題}}
\author{BHB26052 河村太星}
\date{\today}

\begin{document}

\maketitle

講義ノートp.77の記述のように、磁束密度の変動成分$\bm{B}_1$が速度変動$\bm{v}_1$と同様の波動方程式を満たすことを示す。

非圧縮性磁気流体の線形化された方程式のうち、基礎方程式は以下の通りである。
\begin{align}
    \rho_{0}\frac{\partial \bm{v}_{1}}{\partial t} &= \frac{1}{\mu}(\bm{B}_{0}\cdot\nabla)\bm{B}_1 \label{eq:eom} \\
    \frac{\partial \bm{B}_{1}}{\partial t} &= (\bm{B}_{0}\cdot\nabla)\bm{v}_{1} \label{eq:ind}
\end{align}
ここで、$\rho_0$は一様な密度、$\mu$は透磁率、$\bm{B}_0$は背景磁場、$\bm{v}_1$と$\bm{B}_1$はそれぞれ速度および磁束密度の微小変動を表す。
基礎方程式(\ref{eq:ind})の両辺を時間$t$でさらに偏微分すると、以下の関係が得られる。
\begin{equation}
    \frac{\partial^{2}\bm{B}_{1}}{\partial t^{2}} = (\bm{B}_{0}\cdot\nabla)\left(\frac{\partial \bm{v}_{1}}{\partial t}\right)
\end{equation}
この右辺の$\frac{\partial \bm{v}_{1}}{\partial t}$に、運動方程式(\ref{eq:eom})を変形した式
\begin{equation}
    \frac{\partial \bm{v}_{1}}{\partial t} = \frac{1}{\rho_{0}\mu}(\bm{B}_{0}\cdot\nabla)\bm{B}_1
\end{equation}
を代入し、$\rho_0$および$\mu$が空間的に一様な定数であることに注意すると、次式が得られる。
\begin{equation}
    \frac{\partial^{2}\bm{B}_{1}}{\partial t^{2}} = \frac{1}{\rho_{0}\mu}(\bm{B}_{0}\cdot\nabla)^{2} \bm{B}_1
\end{equation}
ここで、講義ノートと同様に、背景磁場が$z$軸方向に一様であるとして$\bm{B}_{0} = (0, 0, B_{0})$と仮定すると、空間微分演算子は次のように書き換えられる。
\begin{equation}
    \bm{B}_{0}\cdot\nabla = B_{0}\frac{\partial}{\partial z}
\end{equation}
これを代入すると、
\begin{align}
    \frac{\partial^{2}\bm{B}_{1}}{\partial t^{2}} &= \frac{1}{\rho_{0}\mu}\left( B_{0}\frac{\partial}{\partial z} \right)\left( B_{0}\frac{\partial}{\partial z} \right) \bm{B}_1 \\
    &= \frac{B_{0}^{2}}{\rho_{0}\mu}\frac{\partial^{2}\bm{B}_{1}}{\partial z^{2}}
\end{align}
両辺を移行して整理すると、最終的に以下の波動方程式を得る。
\begin{equation}
    \frac{\partial^{2}\bm{B}_{1}}{\partial t^{2}} - \frac{B_{0}^{2}}{\mu\rho_{0}}\frac{\partial^{2}\bm{B}_{1}}{\partial z^{2}} = 0 \label{eq:wave}
\end{equation}
式(\ref{eq:wave})は、速度変動$\bm{v}_1$が満たす講義ノートの(6.100)式
\begin{equation}
    \frac{\partial^{2}\bm{v}_{1}}{\partial t^{2}} - \frac{B_{0}^{2}}{\mu\rho_{0}}\frac{\partial^{2}\bm{v}_{1}}{\partial z^{2}} = 0
\end{equation}
と同じ形をしている。したがって、磁束密度の変動$\bm{B}_1$も$\bm{v}_1$と同様に、アルフベン速度$V_A = \sqrt{\frac{B_0^2}{\mu\rho_0}}$で伝播する波動方程式を満たすことが示された。

\end{document}